\section{Conclusion}
\label{sec:conclusion}

Log-sum-exp objectives make responsibilities appear as gradients. This gives
standard neural objectives a precise implicit-EM interpretation, but the
interpretation has a boundary. It is a statement about the coordinates seen by
an LSE or softmax objective, not a guarantee that a whole composed network
inherits EM's optimizer geometry.

The theory identifies that boundary. In distance coordinates, the local LSE
Hessian is uniformly bounded by $1/2$. When component parameters are private,
this yields a parameter-independent curvature bound and explains the
single-site optimizer signature observed in earlier work. When the same site is
trained through a learned dense corridor, the composed CE path contains a
conjugated softmax Hessian whose scale depends on the learned map. The local
guarantee is therefore forfeited.

The supervised experiments confirm the split. Volume control transfers: it
repairs dead units, collapse, and redundancy at an intermediate EM site. The
conditioning does not transfer through ordinary supervised composition: SGD
becomes learning-rate sensitive, Adam becomes useful, and the CE path dominates
the local EM path at the standard auxiliary weight. Cutting that path with
stop-gradient restores the conditioning signature. Increasing the EM weight
until the EM path dominates restores it as well, even with the graph fully
connected.

The resulting rule is local, graded, and structural. Local volume-control
failures can be repaired locally. EM conditioning survives only when
responsibility gradients reach the relevant parameters directly, through a
compatible transport, or with enough dominance to set the effective landscape.
Ordinary dense backpropagation should not be expected to preserve it. This is
the boundary the implicit-EM framework needs: knowing which properties transfer,
which fail, and which architectural interventions can make them survive.
